% !PRINTED 121
\weberpage{121}

\begin{center}
{\large \S.\ 38.}

\vspace{0.5em}
{\large\gesperrt{Beweis des Fundamentalsatzes.}}
\end{center}

Wir gehen nun an den Beweis des \gesperrt{Fundamentalsatzes der Algebra}, dass jede Gleichung $f(x)=0$ wenigstens eine Wurzel hat. Wir schicken einige allgemeine Sätze voraus:

1. \gesperrt{Wenn $S$ ein beliebiges System reeller Zahlen bedeutet, die alle grösser sind als eine bestimmte positive oder negative Zahl $C$ (d. h. ein System, das keine unendlichen negativen Zahlen enthält), so existirt eine untere Grenze für die Zahlen $S$.}

Unter einer \gesperrt{unteren Grenze} $g$ ist eine solche Zahl zu verstehen, die von keiner Zahl des Systems $S$ unterschritten wird, aber so beschaffen, dass, wenn $\delta$ eine beliebig gegebene positive Grösse ist, zwischen $g$ und $g+\delta$ (mit Einschluss der Grenzen) immer noch wenigstens eine Zahl des Systems $S$ liegt, wie klein auch $\delta$ sein mag.

Der Beweis ergiebt sich unmittelbar aus der Möglichkeit eines Schnittes $(\mathfrak A,\mathfrak B)$, den man so construirt, dass man eine Zahl $\alpha$ nach $\mathfrak A$ wirft, wenn sie von keiner Zahl des Systems $S$ unterschritten wird, und eine Zahl $\beta$ nach $\mathfrak B$, wenn sie wenigstens von einer Zahl in $S$ unterschritten wird. Die durch diesen Schnitt bestimmte Zahl $g$ wird von keiner Zahl in $S$ unterschritten, denn sonst gäbe es auch Zahlen, die kleiner als $g$ sind, und also zu $\mathfrak A$ gehören, und die doch von Zahlen in $S$ unterschritten werden, während andererseits jede noch so wenig über $g$ liegende Zahl $\beta$ zu $\mathfrak B$ gehört und also von einer Zahl in $S$ unterschritten wird. Das sind aber die charakteristischen Merkmale der \gesperrt{unteren Grenze}.

2. \gesperrt{Ist $S'$ ein Theil von $S$, so haben auch die Zahlen $S'$ eine untere Grenze $g'$, und diese ist entweder gleich $g$ oder grösser als $g$.}

Denn kleiner als $g$ kann sie nicht sein, da sonst Zahlen in $S'$ und folglich auch in $S$ vorkämen, die unter $g$ liegen.

Es sei nun $f(x)$ eine reelle Function von $x$, die in dem Intervall
\begin{equation}
 a\le x\le b
\tag{1}
\end{equation}
%CONT>>
